\documentclass[a4paper, serif, 10pt]{moderncv}

%% moderncv
% color and style
\moderncvcolor{black}
\moderncvstyle{banking}

%% page layout/margins
\usepackage[top=2.5cm, bottom=2.5cm, left=1.5cm, right=1.5cm]{geometry}

\nopagenumbers{}

%% Babel config for German language
% ref:
% - https://latex3.github.io/babel/guides/locale-german.html
\usepackage[ngerman]{babel}

%% fontspec
\usepackage{fontspec}

\IfFontExistsTF{Linux Libertine}{
  \setmainfont[Ligatures={TeX, Common}, Numbers=OldStyle]{Linux Libertine}
}{}
\IfFontExistsTF{Linux Libertine O}{
  \setmainfont[Ligatures={TeX, Common}, Numbers=OldStyle]{Linux Libertine O}
}{}

%% CTeX
% CJK support, used to show CJK characters in the resume
%
% - heading=false: disable ctex heading
% - punct=kaiming: use kaiming punctuations styles for CJK
% - scheme=plain: use plain scheme, do not override \normalsize font size
% - space=auto: space settings for CJK characters
%
% CJK fonts are resolved per locale language, e.g. Noto Serif/Sans CJK JP
% for Japanese resumes, with fallback to the next region if not installed
%
% ref:
% - http://ctan.mirrorcatalogs.com/language/chinese/ctex/ctex.pdf
\usepackage[UTF8, heading=false, punct=kaiming, scheme=plain, space=auto]{ctex}

\IfFontExistsTF{Noto Serif CJK SC}{
  \setCJKmainfont{Noto Serif CJK SC}
}{}
\IfFontExistsTF{Noto Sans CJK SC}{
  \setCJKsansfont{Noto Sans CJK SC}
}{}

%% line spacing
\usepackage{setspace}
\setstretch{1.125}

%% URL styling - use normal text instead of monospace
\urlstyle{same}

%% Patch \httplink and \httpslink to handle full URLs
% The original moderncv macros always prepend http:// or https:// to URLs.
% This causes issues when the URL already has a protocol (e.g., https://example.com)
% resulting in malformed URLs like https://https://example.com.
% This patch checks if the URL already contains "://" and uses it directly if so.
% Note: We use \str_set:Nx (x-expansion) to expand macros like \@homepage first.
\ExplSyntaxOn
\str_new:N \g_yamlresume_url_str

\RenewDocumentCommand{\httpslink}{O{}m}{
  \str_set:Nx \g_yamlresume_url_str {#2}
  \str_if_in:NnTF \g_yamlresume_url_str {://}
    { \url{#2} }
    { \url{https://#2} }
}

\RenewDocumentCommand{\httplink}{O{}m}{
  \str_set:Nx \g_yamlresume_url_str {#2}
  \str_if_in:NnTF \g_yamlresume_url_str {://}
    { \url{#2} }
    { \url{http://#2} }
}
\ExplSyntaxOff

\name{Dr. Katharina Vogel}{}
\title{Senior HR Managerin – Personalentwicklung, Arbeitsrecht \& Change}
\phone[mobile]{+49 30 5487 2190}
\email{katharina.vogel@hr-vogel.de}
\homepage{https://www.katharina-vogel-hr.de}

\address{Friedrichstraße 112, Berlin, Berlin, Deutschland, 10117}{}{}

\extrainfo{{\small \faLinkedin}\ \href{https://www.linkedin.com/in/katharina-vogel-hr}{@katharina-vogel-hr} {} {} {} • {} {} {} 
{\small \faMedium}\ \href{https://medium.com/@katharina.vogel.hr}{@katharina.vogel.hr}}

\begin{document}

\maketitle

\section{Basisinformationen}

\cvline{}{Senior HR Managerin mit über zehn Jahren Erfahrung in der Personalarbeit mittelständischer und wachsender Technologieunternehmen. Schwerpunkte: strategische Personalentwicklung, Arbeitsrecht, Employer Branding und die Begleitung von Veränderungsprozessen.
%
\begin{itemize}
\item Erfolgreiche Leitung eines HR-Teams von bis zu acht Personen und Verantwortung für rund 450 Mitarbeitende
\item Aufbau datenbasierter HR-Prozesse, Einführung von SAP SuccessFactors und Personio
\item Enge Zusammenarbeit mit Geschäftsführung, Betriebsrat und Führungskräften zur Stärkung einer werteorientierten Unternehmenskultur
\item Ausgeprägte Kommunikationsstärke, diplomatisches Geschick und hohe Umsetzungsorientierung
\end{itemize}}
\section{Ausbildung}

\cventry{Okt. 2010–Sept. 2013}
        {Master, Personalmanagement und Arbeitsrecht, Punkte: 1,7}
        {Humboldt-Universität zu Berlin}
        {\url{https://www.hu-berlin.de}}
        {}
        {\begin{itemize}
\item Schwerpunkt auf strategischer Personalarbeit und arbeitsrechtlichen Fragestellungen
\item Praxisprojekte zu Talentmanagement und Mitarbeiterbindung
\item Ausgezeichnete Studienleistungen und stipendierte Forschungsarbeit
\end{itemize}
\textbf{Kurse}: Arbeitsrecht im Personalwesen, Strategisches Personalmanagement, Organisationspsychologie, Vergütungsmanagement, Change Management, Arbeitsmarktpolitik}

\cventry{Okt. 2007–Sept. 2010}
        {Bachelor, Betriebswirtschaftslehre, Punkte: 1,9}
        {Universität Hamburg}
        {\url{https://www.uni-hamburg.de}}
        {}
        {\begin{itemize}
\item Fundierte betriebswirtschaftliche Grundlagen mit personalwirtschaftlichem Schwerpunkt
\item Teilnahme an Fallstudien zu Personalstrategie und Organisationsentwicklung
\end{itemize}
\textbf{Kurse}: Personalwirtschaft, Organisation und Führung, Wirtschaftsrecht, Statistik, Kosten- und Leistungsrechnung}
\section{Berufserfahrung}

\cventry{Apr. 2021–Heute}
        {Senior HR Managerin / Personalleiterin}
        {Nordlicht Technologies GmbH}
        {\url{https://www.nordlicht-technologies.de}}
        {}
        {\begin{itemize}
\item Strategische und operative Leitung der Personalabteilung für rund 450 Mitarbeitende an drei Standorten
\item Weiterentwicklung der Personalstrategie entlang von Unternehmenswachstum, Internationalisierung und hybriden Arbeitsmodellen
\item Führung und Entwicklung eines Teams aus HR Business Partnern, Recruiting und Personaladministration
\item Verantwortung für arbeitsrechtliche Fragestellungen, Betriebsratszusammenarbeit und Compliance im Personalbereich
\item Einführung von SAP SuccessFactors und Personio zur Professionalisierung von Personalprozessen und HR-Reporting
\item Aufbau eines Employer-Branding-Programms zur Stärkung der Arbeitgeberattraktivität im Technologieumfeld
\end{itemize}
\textbf{Schlüsselwörter}: Personalstrategie, Arbeitsrecht, Change Management, Employer Branding, Führungskräfteentwicklung}

\cventry{Aug. 2017–März 2021}
        {HR Managerin}
        {Berlin Digital Services AG}
        {\url{https://www.berlin-digital-services.de}}
        {}
        {\begin{itemize}
\item Betreuung von Führungskräften und Mitarbeitenden in allen personalrelevanten Themen über den gesamten Employee Lifecycle
\item Verantwortung für Recruiting, Onboarding und Personalentwicklung in einem dynamischen Digitalumfeld
\item Mitwirkung an der Neuausrichtung von Vergütungsstrukturen und Leistungsbeurteilungsprozessen
\item Zusammenarbeit mit dem Betriebsrat bei organisatorischen Veränderungen und personellen Einzelmaßnahmen
\item Einführung standardisierter HR-Kennzahlen und regelmäßiger Personalberichte für die Geschäftsführung
\end{itemize}
\textbf{Schlüsselwörter}: Recruiting, Onboarding, Personalentwicklung, Betriebsrat, HR-Reporting}

\cventry{Juni 2014–Juli 2017}
        {HR Business Partner}
        {Hanseatic Consulting GmbH}
        {\url{https://www.hanseatic-consulting.de}}
        {}
        {\begin{itemize}
\item Beratung von Fachbereichen zu Personalplanung, Talentmanagement und organisatorischen Fragestellungen
\item Begleitung von Reorganisationsprojekten und Entwicklung passender Personalmaßnahmen
\item Unterstützung bei Arbeitsvertragsgestaltung, Zielvereinbarungen und Konfliktlösungen
\item Mitaufbau eines strukturierten Onboarding- und Mentoring-Programms
\end{itemize}
\textbf{Schlüsselwörter}: Talentmanagement, Vergütung, Arbeitsverträge, Konfliktmanagement, Projektarbeit}
\section{Sprachen}

\cvline{Deutsch}{Muttersprache oder zweisprachig \hfill \textbf{Schlüsselwörter}: Muttersprache}
\cvline{Englisch}{Verhandlungssichere Kenntnisse \hfill \textbf{Schlüsselwörter}: Business English, C1}
\cvline{Französisch}{Gute Kenntnisse \hfill \textbf{Schlüsselwörter}: B1, DELF}
\section{Fähigkeiten}

\cvline{Recruiting \& Employer Branding}{Experte \hfill \textbf{Schlüsselwörter}: Talent Acquisition, Active Sourcing, Interviewführung, Candidate Experience, Onboarding}
\cvline{Arbeitsrecht \& Compliance}{Erfahren \hfill \textbf{Schlüsselwörter}: Betriebsverfassungsrecht, Arbeitsvertragsgestaltung, Kündigungsrecht, Personalaktenführung, Datenschutz}
\cvline{Personalentwicklung \& Führung}{Erfahren \hfill \textbf{Schlüsselwörter}: Führungskräftecoaching, Mitarbeitergespräche, Talentmanagement, Nachfolgeplanung, Weiterbildungsplanung}
\cvline{HR-Systeme \& Analytics}{Erfahren \hfill \textbf{Schlüsselwörter}: SAP SuccessFactors, Personio, HR-Kennzahlen, Power BI, Reporting}
\cvline{Change Management \& Kommunikation}{Erfahren \hfill \textbf{Schlüsselwörter}: Change-Kommunikation, Workshop-Moderation, Stakeholder-Management, Konfliktmanagement, Betriebsratskommunikation}
\section{Auszeichnungen}

\cventry{Nov. 2020}
        {Berlin Digital Services AG}
        {HR Excellence Award}
        {}
        {}
        {\begin{itemize}
\item Auszeichnung für die erfolgreiche Einführung eines integrierten Onboarding-Programms
\item Würdigung der messbaren Verbesserung von Time-to-Hire und Mitarbeiterbindung
\end{itemize}}

\cventry{Mai 2022}
        {Nordlicht Technologies GmbH}
        {Top-Arbeitgeber-Zertifizierung}
        {}
        {}
        {\begin{itemize}
\item Mitverantwortung für die Zertifizierung als attraktiver Arbeitgeber im Technologieumfeld
\item Besondere Anerkennung für Diversität, Entwicklungsperspektiven und flexible Arbeitsmodelle
\end{itemize}}
\section{Zertifikate}

\cventry{Juni 2019}
        {Society for Human Resource Management}
        {SHRM Certified Professional (SHRM-CP)}
        {\url{https://www.shrm.org/certification}}
        {}
        {}

\cventry{Sept. 2021}
        {Haufe Akademie}
        {Fachzertifikat Arbeitsrecht im Personalwesen}
        {\url{https://www.haufe-akademie.de}}
        {}
        {}

\cventry{März 2022}
        {Hasso-Plattner-Institut}
        {Design Thinking Facilitator}
        {\url{https://hpi.de}}
        {}
        {}
\section{Veröffentlichungen}

\cventry{Sept. 2022}
        {Hybrides Arbeiten: Leitfaden für moderne Personalarbeit}
        {Personalwirtschaft Verlag}
        {\url{https://www.personalwirtschaft.de}}
        {}
        {\begin{itemize}
\item Praxisnaher Leitfaden zur Gestaltung hybrider Arbeitsmodelle
\item Handlungsempfehlungen für Führung, Zusammenarbeit und Mitarbeiterbindung
\item Fallbeispiele aus mittelständischen Technologieunternehmen
\end{itemize}}
\section{Referenzen}

\cventry{\emaillink[martin.seidel@nordlicht-technologies.de]{martin.seidel@nordlicht-technologies.de}}
        {Geschäftsführer, Nordlicht Technologies GmbH}
        {Dr. Martin Seidel}
        {+49 30 5487 2100}
        {}
        {Frau Dr. Vogel hat unsere Personalarbeit strategisch neu ausgerichtet und durch klare Prozesse, transparente Kommunikation und eine starke Führungskräfteentwicklung nachhaltig geprägt. Sie verbindet arbeitsrechtliche Sicherheit mit unternehmerischem Denken.}
\section{Projekte}

\cventry{März 2019–Dez. 2023}
        {Ein unternehmensweites Mentoring-Programm zur gezielten Förderung von Frauen in Führungspositionen.}
        {Mentoring-Programm Female Leaders}
        {\url{https://www.nordlicht-technologies.de/karriere/female-leaders}}
        {}
        {\begin{itemize}
\item Konzeption und Steuerung eines strukturierten Mentoring-Programms mit externen Coachings
\item Aufbau von Tandempartnerschaften zwischen Nachwuchsführungskräften und erfahrenen Führungspersönlichkeiten
\item Regelmäßige Erfolgsmessung über Teilnehmerfeedback, Beförderungsquoten und Retention
\item Aktive Einbindung der Geschäftsführung und des Betriebsrats
\end{itemize}
\textbf{Schlüsselwörter}: Mentoring, Diversity, Führungskräfteentwicklung}

\cventry{Jan. 2022–Dez. 2022}
        {Aufbau eines Dashboards zur Visualisierung zentraler HR-Kennzahlen für Geschäftsführung und Führungskräfte.}
        {HR-Analytics-Dashboard}
        {}
        {}
        {\begin{itemize}
\item Definition relevanter Kennzahlen zu Fluktuation, Recruiting, Krankheit und Personalentwicklung
\item Entwicklung standardisierter Reports in Power BI auf Basis bereinigter Personaldaten
\item Schulung von Führungskräften zur datenbasierten Personalsteuerung
\item Verbesserung der Entscheidungsgrundlagen im monatlichen Management-Review
\end{itemize}
\textbf{Schlüsselwörter}: HR-Kennzahlen, Power BI, Reporting, Datenqualität}
\section{Interessen}

\cvline{Berufliches Engagement}{Diversity \& Inclusion, Mentoring, Bildungsgerechtigkeit}
\cvline{Ausgleich}{Laufen, Yoga, Wandern}
\section{Ehrenamtliche Tätigkeiten}

\cventry{Jan. 2019–Heute}
        {Mentorin für Geflüchtete}
        {Berliner Netzwerk für berufliche Integration e.V.}
        {\url{https://www.berufliche-integration-berlin.de}}
        {}
        {\begin{itemize}
\item Unterstützung von Geflüchteten bei der beruflichen Orientierung und Bewerbung
\item Durchführung von Bewerbungstrainings und Einzelcoachings
\item Vernetzung mit Unternehmen und Bildungsträgern
\item Mitwirkung an Veranstaltungen zur interkulturellen Öffnung des Arbeitsmarkts
\end{itemize}}
    
\end{document}